% Part:sets-functions-relations
% Chapter: sets
% Section: schroder-bernstein

\documentclass[../../../include/open-logic-section]{subfiles}

\begin{document}

\olfileid{sfr}{siz}{sb}

\olsection{Gagasan Gedhene Himpunan, lan Schr\"oder-Bernstein}

\begin{explain}
Ana gagasan intuisi mangkene: yen $A$ ora luwih gedhe tinimbang $B$ lan $B$
ora luwih gedhe tinimbang $A$, mula $A$ lan $B$ padha cacahe.
Satemene, yen gagasan iki \emph{kleru}, kita bakal angel mbenerake
anggepan yen definisi padha cacahe ana gandhengane karo mbandhingake
"gedhene" himpunan!{} Nanging untunge gagasan intuisi iki bener.
Iki dibuktekake dening Teorema Schr\"oder-Bernstein.
\end{explain}

\begin{thm}[Schr\"oder-Bernstein]
	\ollabel{thm:schroder-bernstein}
	Yen $\cardle{A}{B}$ lan $\cardle{B}{A}$,
	mula $\cardeq{A}{B}$.
\end{thm}

\begin{explain}
Kanthi tembung liya, yen ana !!a{injection} saka $A$ menyang~$B$, lan
!!a{injection} saka $B$ menyang~$A$, mula ana !!a{bijection} saka $A$
menyang~$B$.

Nanging asil iki pancen rada \emph{angel} dibuktekake.
Sanajan Cantor kang ngandharake asil iki, wong liya kang
mbuktekake.\footnote{Babagan sejarahe luwih lanjut, delengen upamane
\citet[kaca~165--6]{Potter2004}.}
\oliflabeldef{sfr:cardinals:card-sb:sec}{Kita lagi bakal bisa
\emph{mbuktekake} Schr\"oder-Bernstein ing
\olref[sfr][cardinals][card-sb]{sec}.}{}%
Saiki kita bisa (lan kudu) nampa asil iki dhisik tanpa bukti.

Untunge, Schr\"oder-Bernstein iku \emph{bener}, lan mbenerake anggepan
yen relasi kang wis ditetepake, yaiku $\cardeq{A}{B}$ lan $\cardle{A}{B}$,
ana gandhengane karo "gedhene". Kajaba iku, Schr\"oder-Bernstein
banget \emph{migunani}. Kadhangkala angel nemokake !!a{bijection}
antarane rong himpunan kang padha cacahe. Teorema Schr\"oder-Bernstein
ngidini kita mecah pambandhingan dadi rong arah, saengga mung perlu
nemokake !!a{injection} saka himpunan kapisan menyang kapindho,
lan kosok baline.
\end{explain}
\end{document}
